%========================================
% MatinBook — Stage 09: Algorithm Test
% Version: v1.1
% Purpose:
%   Validate the algorithm system of MatinBook v1.1:
%     1. Standard algorithm + algorithmic environments
%     2. Persian captions via \caption inside algorithm
%     3. Automatic numbering (chapter.algorithm)
%     4. List of algorithms
%     5. Cross-references with \cref
%     6. Persian text inside pseudocode comments
%
% Compilation:
%   xelatex -shell-escape stage09-algorithms.tex
%   xelatex -shell-escape stage09-algorithms.tex
%   (2 passes required for cover + cross-references)
%
% Expected outcome:
%   A professionally typeset PDF demonstrating that every
%   algorithm renders correctly, with a Persian caption, a
%   continuous chapter-based number, and working references.
%========================================

%------------------------------------------------------------------------
% Search Path (for tests directory)
%------------------------------------------------------------------------
\makeatletter
\def\input@path{{../}}
\makeatother

%------------------------------------------------------------------------
% Document Class
%------------------------------------------------------------------------
\documentclass{matinbook}

%------------------------------------------------------------------------
% Book Metadata
%------------------------------------------------------------------------
\title{آزمون سیستم الگوریتم‌ها}
\author{تیم توسعه MatinBook}
\date{مهر ۱۴۰۵}

\booktitle{آزمون سیستم الگوریتم‌ها}

\renewcommand{\repository}{github.com/matinmahmoudi-cs/matinbook}

%========================================================================
% DOCUMENT
%========================================================================
\begin{document}

%------------------------------------------------------------------------
% FRONT COVER
%------------------------------------------------------------------------
\makecover
    {آزمون سیستم الگوریتم‌ها}
    {ارزیابی شبه‌کد، کپشن و ارجاع در MatinBook}
    {تیم توسعه MatinBook}
    {مهر ۱۴۰۵}

%------------------------------------------------------------------------
% FRONT MATTER (symmetric margins)
%------------------------------------------------------------------------
\frontmatter
\frontmattergeometry

% Title page
\maketitle

% Copyright / Colophon
\thispagestyle{empty}
\vfill
\begin{center}
    {\large\textbf{شناسنامه آزمون}}

    \vspace{0.8cm}

    \begin{tabular}{rl}
        عنوان: & آزمون سیستم الگوریتم‌ها \\
        نسخه: & \matinbookversion \\
        تاریخ: & \matinbookdate \\
        موتور: & \lr{XeLaTeX} \\
        مجوز: & \lr{MIT} \\
    \end{tabular}

    \vspace{0.8cm}

    {\small این سند بخشی از مجموعه آزمون‌های یکپارچگی \lr{MatinBook} است.}
\end{center}
\clearpage

% Preface
\chapter*{پیشگفتار}
\addcontentsline{toc}{chapter}{پیشگفتار}

این سند، نهمین آزمون از مجموعه‌ی آزمون‌های یکپارچگی
\lr{MatinBook} نسخه‌ی \lr{\matinbookversion} است. هدف آن،
اعتبارسنجی کامل سیستم الگوریتم‌های کلاس در شش حوزه‌ی
متمایز است: محیط‌های استاندارد \lr{algorithm} و
\lr{algorithmic}، کپشن‌های فارسی، شماره‌گذاری خودکار،
فهرست الگوریتم‌ها، ارجاع هوشمند و متن فارسی داخل
شبه‌کد.

در کتاب‌های برنامه‌نویسی و علوم کامپیوتر، نمایش
حرفه‌ای الگوریتم‌ها اهمیت ویژه‌ای دارد. هر الگوریتم
باید با شماره‌ی منظم، کپشن گویا و قابلیت ارجاع نمایش
داده شود. \lr{MatinBook} با استفاده از بسته‌ی
\lr{algorithm} و یکپارچگی با \lr{xepersian}، این
اهداف را محقق می‌سازد.

هر بخش از این آزمون، با مثال‌های عملی و ارجاعات متقابل
همراه است تا خواننده بتواند درستی رفتار هر زیرسیستم را
در بافت واقعی یک کتاب ارزیابی کند.

\clearpage

% Table of Contents
\tableofcontents
\clearpage

% List of Figures
\listoffigures
\clearpage

% List of Tables
\listoftables
\clearpage

% List of Algorithms
\listofalgorithms
\clearpage

%------------------------------------------------------------------------
% MAIN MATTER (wide margin for notes)
%------------------------------------------------------------------------
\mainmatter
\mainmattergeometry

%========================================================================
% CHAPTER 1 — SEARCH ALGORITHMS
%========================================================================
\chapter{الگوریتم‌های جستجو}
\label{chap:search}

\section{مقدمه}
\label{sec:search-intro}

الگوریتم‌های جستجو\index{الگوریتم!جستجو} از پایه‌ای‌ترین
الگوریتم‌های علوم کامپیوتر هستند. در این فصل، دو
الگوریتم اصلی جستجو را بررسی می‌کنیم.

\mnote{الگوریتم جستجو، یک مقدار مشخص را در یک ساختار داده پیدا می‌کند.}

\section{جستجوی خطی}
\label{sec:search-linear}

ساده‌ترین روش جستجو، بررسی تک‌تک عناصر آرایه است.
این روش برای آرایه‌های مرتب و نامرتب کار می‌کند.
الگوریتم \cref{alg:linear-search} این روش را نشان می‌دهد:

\begin{latin}
\begin{algorithm}
\caption{جستجوی خطی}
\label{alg:linear-search}
\begin{algorithmic}[1]
    \Require $A[1..n]$: array of $n$ elements
    \Require $x$: target value to find
    \Ensure Index $i$ such that $A[i] = x$, or $-1$ if not found

    \For{$i \gets 1$ \textbf{to} $n$}
        \If{$A[i] = x$}
            \State \Return $i$
        \EndIf
    \EndFor
    \State \Return $-1$
\end{algorithmic}
\end{algorithm}
\end{latin}

\begin{theorem}[پیچیدگی جستجوی خطی]
    الگوریتم \cref{alg:linear-search} در بدترین حالت
    \lr{$O(n)$} مقایسه انجام می‌دهد، زیرا ممکن است مجبور
    شود تمام $n$ عنصر آرایه را بررسی کند.
    \label{thm:linear-complexity}
\end{theorem}

\section{جستجوی دودویی}
\label{sec:search-binary}

برای آرایه‌های \textbf{مرتب}، جستجوی دودویی بسیار
کارآمدتر از جستجوی خطی است. الگوریتم
\cref{alg:binary-search} این روش را نشان می‌دهد:

\begin{latin}
\begin{algorithm}
\caption{جستجوی دودویی}
\label{alg:binary-search}
\begin{algorithmic}[1]
    \Require $A[1..n]$: sorted array (ascending order)
    \Require $x$: target value to find
    \Ensure Index $i$ such that $A[i] = x$, or $-1$ if not found

    \State $left \gets 1$
    \State $right \gets n$

    \While{$left \leq right$}
        \State $mid \gets \lfloor (left + right) / 2 \rfloor$
        \If{$A[mid] = x$}
            \State \Return $mid$
        \ElsIf{$A[mid] < x$}
            \State $left \gets mid + 1$
        \Else
            \State $right \gets mid - 1$
        \EndIf
    \EndWhile

    \State \Return $-1$
\end{algorithmic}
\end{algorithm}
\end{latin}

\begin{theorem}[پیچیدگی جستجوی دودویی]
    الگوریتم \cref{alg:binary-search} در بدترین حالت
    \lr{$O(\log n)$} مقایسه انجام می‌دهد، زیرا در هر
    مرحله فضای جستجو را نصف می‌کند.
    \label{thm:binary-complexity}
\end{theorem}

\section{مقایسه دو الگوریتم}
\label{sec:search-comparison}

مقایسه‌ی \cref{alg:linear-search} و
\cref{alg:binary-search} نشان می‌دهد که برای آرایه‌های
مرتب، جستجوی دودویی به‌مراتب سریع‌تر است. جدول
\cref{tab:search-comparison} این مقایسه را خلاصه می‌کند:

\begin{matintable}{مقایسه الگوریتم‌های جستجو}{tab:search-comparison}
\begin{tblr}{
    width = \textwidth,
    colspec = {l X[c] X[c] X[c]},
    hlines, vlines,
    colsep = 8pt,
    rowsep = 4pt,
    row{1} = {font=\bfseries},
}
معیار & جستجوی خطی & جستجوی دودویی & برنده \\
پیچیدگی زمانی & \lr{$O(n)$} & \lr{$O(\log n)$} & دودویی \\
نیاز به مرتب‌سازی & خیر & \textcolor{matinred}{بله} & خطی \\
پیاده‌سازی & ساده & متوسط & خطی \\
کارایی برای $n$ کوچک & خوب & سربار بیشتر & خطی \\
کارایی برای $n$ بزرگ & ضعیف & \textcolor{matingreen}{عالی} & دودویی \\
\end{tblr}
\end{matintable}

%========================================================================
% CHAPTER 2 — SORTING ALGORITHMS
%========================================================================
\chapter{الگوریتم‌های مرتب‌سازی}
\label{chap:sorting}

\section{مرتب‌سازی سریع}
\label{sec:sorting-quicksort}

\lr{Quick Sort}\index{الگوریتم!مرتب‌سازی سریع} یکی از
سریع‌ترین الگوریتم‌های مرتب‌سازی در عمل است. الگوریتم
\cref{alg:quicksort} این روش را نشان می‌دهد:

\begin{latin}
\begin{algorithm}
\caption{مرتب‌سازی سریع (\lr{Quick Sort})}
\label{alg:quicksort}
\begin{algorithmic}[1]
    \Require $A[1..n]$: array of $n$ elements
    \Require $low, high$: indices defining the subarray
    \Ensure Array $A$ is sorted in ascending order

    \Function{QuickSort}{$A$, $low$, $high$}
        \If{$low < high$}
            \State $p \gets \Call{Partition}{A, low, high}$
            \State $\Call{QuickSort}{A, low, p - 1}$
            \State $\Call{QuickSort}{A, p + 1, high}$
        \EndIf
    \EndFunction
    \Statex
    \Function{Partition}{$A$, $low$, $high$}
        \State $pivot \gets A[high]$
        \State $i \gets low - 1$
        \For{$j \gets low$ \textbf{to} $high - 1$}
            \If{$A[j] \leq pivot$}
                \State $i \gets i + 1$
                \State Swap $A[i]$ and $A[j]$
            \EndIf
        \EndFor
        \State Swap $A[i+1]$ and $A[high]$
        \State \Return $i + 1$
    \EndFunction
\end{algorithmic}
\end{algorithm}
\end{latin}

\section{الگوریتم اقلیدس}
\label{sec:sorting-euclidean}

الگوریتم اقلیدس\index{الگوریتم!اقلیدس} برای محاسبه‌ی
بزرگ‌ترین مقسوم‌علیه مشترک (\lr{GCD}) در
\cref{alg:euclidean} نمایش داده شده است:

\begin{latin}
\begin{algorithm}
\caption{الگوریتم اقلیدس برای $\gcd(a, b)$}
\label{alg:euclidean}
\begin{algorithmic}[1]
    \Require $a, b$: non-negative integers
    \Ensure $\gcd(a, b)$: greatest common divisor

    \While{$b \neq 0$}
        \State $r \gets a \bmod b$
        \State $a \gets b$
        \State $b \gets r$
    \EndWhile
    \State \Return $a$
\end{algorithmic}
\end{algorithm}
\end{latin}

\begin{example}
    محاسبه‌ی $\gcd(48, 18)$ با الگوریتم
    \cref{alg:euclidean}:
    \begin{align*}
        \gcd(48, 18) &: \; 48 \bmod 18 = 12 \\
        \gcd(18, 12) &: \; 18 \bmod 12 = 6 \\
        \gcd(12, 6)  &: \; 12 \bmod 6 = 0 \\
        \gcd(6, 0)   &: \; \text{نتیجه} = 6
    \end{align*}
    \label{ex:gcd-example}
\end{example}

%========================================================================
% CHAPTER 3 — GRAPH ALGORITHMS
%========================================================================
\chapter{الگوریتم‌های گراف}
\label{chap:graph}

\section{جستجوی اول سطح (BFS)}
\label{sec:graph-bfs}

\lr{Breadth-First Search}\index{الگوریتم!BFS} گراف را
سطح به سطح پیمایش می‌کند. الگوریتم \cref{alg:bfs}
این روش را نشان می‌دهد:

\begin{latin}
\begin{algorithm}
\caption{جستجوی اول سطح (\lr{BFS})}
\label{alg:bfs}
\begin{algorithmic}[1]
    \Require $G = (V, E)$: an unweighted graph
    \Require $s$: source vertex
    \Ensure $dist[v]$: shortest distance from $s$ to each $v \in V$

    \ForAll{$v \in V$}
        \State $dist[v] \gets \infty$
        \State $visited[v] \gets \text{false}$
    \EndFor

    \State $dist[s] \gets 0$
    \State $visited[s] \gets \text{true}$
    \State Create an empty queue $Q$
    \State $Q.\text{enqueue}(s)$

    \While{$Q$ is not empty}
        \State $u \gets Q.\text{dequeue}()$
        \ForAll{neighbor $v$ of $u$}
            \If{$\neg visited[v]$}
                \State $visited[v] \gets \text{true}$
                \State $dist[v] \gets dist[u] + 1$
                \State $Q.\text{enqueue}(v)$
            \EndIf
        \EndFor
    \EndWhile

    \State \Return $dist$
\end{algorithmic}
\end{algorithm}
\end{latin}

\section{الگوریتم دایکسترا}
\label{sec:graph-dijkstra}

الگوریتم دایکسترا\index{الگوریتم!دایکسترا} کوتاه‌ترین
مسیر در گراف وزن‌دار را پیدا می‌کند. الگوریتم
\cref{alg:dijkstra} این روش را نشان می‌دهد:

\begin{latin}
\begin{algorithm}
\caption{الگوریتم دایکسترا (\lr{Dijkstra's Algorithm})}
\label{alg:dijkstra}
\begin{algorithmic}[1]
    \Require $G = (V, E)$: weighted graph with non-negative weights
    \Require $s$: source vertex
    \Ensure $dist[v]$: shortest distance from $s$ to each $v \in V$

    \ForAll{$v \in V$}
        \State $dist[v] \gets \infty$
        \State $prev[v] \gets \text{null}$
    \EndFor
    \State $dist[s] \gets 0$
    \State $PQ \gets$ priority queue of vertices keyed by $dist$
    \State Insert all vertices into $PQ$

    \While{$PQ$ is not empty}
        \State $u \gets$ extract vertex with minimum $dist$ from $PQ$
        \ForAll{neighbor $v$ of $u$}
            \State $alt \gets dist[u] + weight(u, v)$
            \If{$alt < dist[v]$}
                \State $dist[v] \gets alt$
                \State $prev[v] \gets u$
                \State Decrease key of $v$ in $PQ$ to $alt$
            \EndIf
        \EndFor
    \EndWhile

    \State \Return $(dist, prev)$
\end{algorithmic}
\end{algorithm}
\end{latin}

%========================================================================
% CHAPTER 4 — ALGORITHM ANALYSIS
%========================================================================
\chapter{تحلیل و مقایسه‌ی الگوریتم‌ها}
\label{chap:analysis}

\section{ارجاع به الگوریتم‌ها}
\label{sec:analysis-refs}

سیستم ارجاع \lr{MatinBook} امکان ارجاع به الگوریتم‌ها
را با استفاده از \lr{\textbackslash cref} فراهم می‌کند:

\begin{itemize}
    \item \cref{alg:linear-search} — ساده‌ترین روش جستجو
    با پیچیدگی \lr{$O(n)$}
    \item \cref{alg:binary-search} — جستجوی سریع با
    پیچیدگی \lr{$O(\log n)$} برای آرایه‌های مرتب
    \item \cref{alg:quicksort} — مرتب‌سازی سریع با
    پیچیدگی متوسط \lr{$O(n \log n)$}
    \item \cref{alg:euclidean} — الگوریتم باستانی اقلیدس
    با پیچیدگی \lr{$O(\log \min(a,b))$}
    \item \cref{alg:bfs} — پیمایش سطح به سطح گراف با
    پیچیدگی \lr{$O(V + E)$}
    \item \cref{alg:dijkstra} — کوتاه‌ترین مسیر با
    پیچیدگی \lr{$O((V+E) \log V)$}
\end{itemize}

\section{جدول مقایسه}
\label{sec:analysis-comparison}

جدول \cref{tab:algorithm-comparison} پیچیدگی الگوریتم‌های
معرفی‌شده را مقایسه می‌کند:

\begin{matintable}{مقایسه پیچیدگی الگوریتم‌ها}{tab:algorithm-comparison}
\begin{tblr}{
    width = \textwidth,
    colspec = {l X[c] X[c] X[c]},
    hlines, vlines,
    colsep = 8pt,
    rowsep = 4pt,
    row{1} = {font=\bfseries},
}
الگوریتم & کاربرد & پیچیدگی & نیازمندی \\
جستجوی خطی & جستجوی ساده & \lr{$O(n)$} & ندارد \\
جستجوی دودویی & جستجوی سریع & \lr{$O(\log n)$} & آرایه مرتب \\
مرتب‌سازی سریع & مرتب‌سازی & \lr{$O(n \log n)$} & ندارد \\
الگوریتم اقلیدس & محاسبه GCD & \lr{$O(\log n)$} & اعداد صحیح \\
\lr{BFS} & پیمایش گراف & \lr{$O(V+E)$} & گراف بدون وزن \\
دایکسترا & کوتاه‌ترین مسیر & \lr{$O((V+E)\log V)$} & یال‌های غیرمنفی \\
\end{tblr}
\end{matintable}

%========================================================================
% CHAPTER 5 — SUMMARY
%========================================================================
\chapter{خلاصه‌ی نتایج}
\label{chap:summary}

\section{وضعیت سیستم الگوریتم}
\label{sec:summary-status}

جدول \cref{tab:algorithm-status} وضعیت قابلیت‌های
الگوریتم \lr{MatinBook} را نشان می‌دهد.

\begin{matintable}{وضعیت قابلیت‌های الگوریتم}{tab:algorithm-status}
\begin{tblr}{
    width = \textwidth,
    colspec = {l X[c] X[c]},
    hlines, vlines,
    colsep = 8pt,
    rowsep = 4pt,
    row{1} = {font=\bfseries},
}
قابلیت & توضیحات & وضعیت \\
شماره‌گذاری & شماره‌گذاری خودکار «فصل.شماره» & \textcolor{matingreen}{فعال} \\
عنوان فارسی & نمایش عنوان با \lr{\textbackslash caption} & \textcolor{matingreen}{فعال} \\
فهرست الگوریتم‌ها & \lr{\textbackslash listofalgorithms} & \textcolor{matingreen}{فعال} \\
ارجاع & \lr{\textbackslash cref} با پیشوند «الگوریتم» & \textcolor{matingreen}{فعال} \\
\lr{Require/Ensure} & پیش‌نیاز و خروجی الگوریتم & \textcolor{matingreen}{فعال} \\
\lr{For/While/If} & ساختارهای کنترلی استاندارد & \textcolor{matingreen}{فعال} \\
\lr{Function} & تعریف و فراخوانی تابع & \textcolor{matingreen}{فعال} \\
محیط \lr{latin} & شبه‌کد درون محیط لاتین & \textcolor{matingreen}{فعال} \\
\end{tblr}
\end{matintable}

\section{معیارهای ارزیابی}
\label{sec:summary-criteria}

در این آزمون، سیستم الگوریتم \lr{MatinBook} بر اساس
پنج معیار ارزیابی شد:

\begin{enumerate}
    \item \textbf{شماره‌گذاری}: منظم و پیوسته در هر فصل.
    \item \textbf{کپشن}: فارسی و گویا با \lr{\textbackslash caption}.
    \item \textbf{ارجاع}: هوشمند با \lr{\textbackslash cref}.
    \item \textbf{ساختار}: استاندارد \lr{algorithmic}.
    \item \textbf{فهرست}: \lr{\textbackslash listofalgorithms}.
\end{enumerate}

\section{نتیجه‌گیری}
\label{sec:summary-conclusion}

\textbf{تمام زیرسیستم‌های الگوریتم \lr{MatinBook} نسخه‌ی
\lr{\matinbookversion} با موفقیت بارگذاری و آزمایش شدند.
شش الگوریتم مختلف با عناوین فارسی، شماره‌گذاری خودکار،
فهرست الگوریتم‌ها و ارجاع هوشمند به‌درستی کار می‌کنند.}

%------------------------------------------------------------------------
% BACK MATTER (symmetric margins)
%------------------------------------------------------------------------
\backmatter
\frontmattergeometry

% Bibliography (empty placeholder)
\printbibliography[title={منابع و مراجع}]

% Index
\printindex

%------------------------------------------------------------------------
% BACK COVER
%------------------------------------------------------------------------
\authorbio{%
    تیم توسعه‌ی \lr{MatinBook}، گروهی از پژوهشگران و برنامه‌نویسان
    فارسی‌زبان است که با هدف ارائه‌ی یک راه‌حل حرفه‌ای برای
    حروف‌چینی کتاب‌های فنی فارسی فعالیت می‌کند. این پروژه حاصل
    سال‌ها تجربه در حوزه‌ی نشر دیجیتال و برنامه‌نویسی است.
}

\makebackcover{%
    این سند، نهمین آزمون از مجموعه‌ی آزمون‌های یکپارچگی
    \lr{MatinBook} نسخه‌ی \lr{\matinbookversion} است.

    \vspace{0.5cm}

    \textbf{زیرسیستم‌های آزموده‌شده:}
    \begin{itemize}
        \item محیط‌های \lr{algorithm} و \lr{algorithmic}
        \item کپشن‌های فارسی
        \item شماره‌گذاری خودکار
        \item فهرست الگوریتم‌ها
        \item ارجاع هوشمند با \lr{\textbackslash cref}
        \item متن فارسی داخل شبه‌کد
    \end{itemize}
}

\end{document}